\section*{Chapter \ref{ch:nw:switches-and-layer-1-devices} --- Switches and Layer-1 Devices}

\begin{solution}{exo:nw:switches-and-layer-1-devices:1}
Store-and-forward: $8 \times 512 / 25 + 300 = 163.8 + 300 = \qty{463.8}{\nano\second}$. Cut-through:
$8 \times 64 / 25 + 300 = 20.5 + 300 = \qty{320.5}{\nano\second}$.
\end{solution}

\begin{solution}{exo:nw:switches-and-layer-1-devices:2}
From 25 in to 10 out it can cut through: the bits arrive faster than they leave, so the output never runs dry (the difference
accumulates in the buffer until the frame ends). From 10 in to 25 out it cannot: the output would send faster than the bits
arrive and run out in the middle of the frame, so the frame is stored first.
\end{solution}

\begin{solution}{exo:nw:switches-and-layer-1-devices:3}
$2 \times 6 = \qty{12}{\giga\bit\per\second}$ offered to a \qty{10}{\giga\bit\per\second} port: $1 - 10/12 = 1/6$, about 17\% of
the copy is lost while both directions are that busy.
\end{solution}

\begin{solution}{exo:nw:switches-and-layer-1-devices:4}
$D$ has standard deviation $100\sqrt2 = 141.4$~ns: $2\Phi(96/141.4) - 1 = 2\Phi(0.679) - 1 = 0.50$.
\end{solution}

\begin{solution}{exo:nw:switches-and-layer-1-devices:5}
Serialisation $8 \times 104/10 = 83.2$~ns, fibre $36 \times 4.88 = 175.6$~ns, devices 4~ns. The fibre: shorter cables inside the cage
and, above all, a shorter cross-connect, which is the venue's to decide (chapter~9); the serialisation next, by a faster link.
\end{solution}

\begin{solution}{exo:nw:switches-and-layer-1-devices:6}
A fan-out copies one input to outputs that carry nothing else: an output is never busy when its bits arrive, so nothing waits. A
multiplexer's output is shared by several inputs, and a frame arriving while another is being sent must wait for it. The fan-out's
latency is a constant; the multiplexer's is a constant plus a queueing delay that depends on the other inputs, its jitter.
\end{solution}

\begin{solution}{exo:nw:switches-and-layer-1-devices:7}
Each multiplexer now serves two servers: the mean wait at \qty{50}{\nano\second} falls from 93.2~ns (four on one) to 22.7~ns (two
on each), the value of \texttt{mux\_contention(2, 50)}; the design gains a handoff and a multiplexer.
\end{solution}

\begin{solution}{exo:nw:switches-and-layer-1-devices:8}
One box has one power path, one chassis, one set of optics and one configuration: a power supply or a configuration change takes
both lines at once. Put line A and line B on different devices, as the method says; lack of software removes one failure mode,
not the others.
\end{solution}

\begin{solution}{pb:nw:switches-and-layer-1-devices:1}
\begin{enumerate}
\item Commodity $1\,439.8 + 1\,430.2 = 2\,870.0$~ns; cut-through $560.0 + 556.8 = 1\,116.7$~ns; layer-1 $262.8 + 294.6 = 557.3$~ns.
\item 1\,753.3~ns; then 559.4~ns.
\item Two serialisations of the 104-byte frame, $2 \times 83.2 = 166.4$~ns, plus \qty{1}{\micro\second} of the model's internal time.
\item $36 \times 4.88 = 175.6$~ns.
\item 93.2~ns.
\item 186.1~ns.
\item Yes: $294.6 + 186.1 = 480.7$~ns against the cut-through order path's 556.8~ns.
\item Across the two multiplexers, two each, with the two most correlated strategies on different ones.
\item Both switches, the cable between them and the two cables to the servers.
\item Each line has its own switch and each switch reaches both servers and its own order handoff: losing one switch loses one
  line and one order path, never all of either.
\item Two order handoffs, each on its own multiplexer, and the venue accepting orders on both (two sessions).
\item Nothing for the servers: they take the other line, which has the same latency; the firm loses its protection against the next
  failure until the device is replaced.
\item $60\,000 / 1\,753.3 \approx 34.2$ per nanosecond a year.
\item $90\,000 / 559.4 \approx 160.9$ per nanosecond a year.
\item \textbf{Named result.} Trigger paths of 2\,870, 1\,117 and 557~ns; the first 1\,753 nanoseconds cost about 34 a year each, the
  next 559 about 161 each.
\item One that wins or loses races decided by hundreds of nanoseconds on this venue, such as the latency-arbitrage and market-making
  strategies of One Quant Book~11.
\item Filtering by address or group, routing between segments, per-port queues and counters of a switch; the design must do its own
  filtering on the server.
\item From the cables (the cross-connect's length), from serialisation (faster links) and from the servers themselves (chapters~3
  and~8, and One Quant Book~13).
\item The single-point check, the tap check, and a capture of the real paths against the computed ones (chapter~5).
\item One that does not read the frame at all.
\end{enumerate}
\end{solution}

\begin{solution}{iq:nw:switches-and-layer-1-devices:1}
Store-and-forward receives the whole frame, checks it and then sends it: it adds a full serialisation. Cut-through starts sending
after the header: it adds a constant. It stores and forwards when the output is busy or faster than the input, and on switches
that must rewrite the frame.
\iqlookfor{the serialisation term and the fallback conditions.}
\end{solution}

\begin{solution}{iq:nw:switches-and-layer-1-devices:2}
A device that connects ports at the signal level, regenerating bits without decoding frames: nanoseconds, no queues, fan-out and
patching under software control. It cannot read addresses, filter groups, route, or merge traffic from several ports (a
multiplexer, which is a different device, does that with a queue).
\iqlookfor{no framing, hence no filtering and no contention.}
\end{solution}

\begin{solution}{iq:nw:switches-and-layer-1-devices:3}
Orders go many to one; a fan-out goes one to many. A multiplexer merges them with a small queue: when several servers answer the same
event within a frame's serialisation, their orders queue behind one another, and a firm's fastest strategies collide exactly when
they matter. Spread correlated strategies over several handoffs.
\iqlookfor{contention that correlates with market events.}
\end{solution}

\begin{solution}{iq:nw:switches-and-layer-1-devices:4}
A tap: it is on the link, copies every bit including errors at the same time as the original, cannot drop the copy, and keeps
working when the switch does not. A mirror port is a switch output that can be oversubscribed and whose copy passes through the
switch's own queueing, so its timing is not the wire's.
\iqlookfor{oversubscription and timing fidelity of the copy.}
\end{solution}

\begin{solution}{iq:nw:switches-and-layer-1-devices:5}
Line A and line B through separate taps to separate fan-outs, each fanning out to all four servers; orders from each server through
a multiplexer to the handoff, and a second handoff and multiplexer if the venue allows; capture on every tap; management out of band.
Failures: a fan-out (servers use the other line), a multiplexer (servers use the second handoff, or stop if there is none: the
single point to negotiate away), a cable, a server (its strategies move).
\iqlookfor{no shared device between A and B, and an honest statement of what remains single.}
\end{solution}

\begin{solution}{iq:nw:switches-and-layer-1-devices:6}
At the open the output ports are busy: frames that find their output sending another frame are stored and wait, so the switch adds
queueing (and loses cut-through for those frames). Look at the \omterm{def:nw:networking-for-trading:egress}{egress queue} occupancy and at \omterm{def:nw:networking-for-trading:burst}{microbursts} (chapter~1), not at the
data sheet's idle latency.
\iqlookfor{load-dependent queueing and the fallback to store-and-forward.}
\end{solution}
